% GENERATED FILE. Edit canonical lessons, then run npm run generate:books.
% canonical-source-hash: fnv1a64:f2a718244bcff433
% canonical-lessons: TE-C279-koluko, TE-C279-ubbu, TE-C279-kallu-tirugu, TE-C279-taku, TE-C279-olucu

\chapter{To recover, To swell, To feel dizzy, To touch, To peel}
\label{ch:te-to-recover-to-swell-to-feel-dizzy-to-touch-to-peel}
\hlchaptermodality{279}

\begin{chapteropening}
\textbf{By the end of this chapter:} \emph{I can say to recover, to swell, to feel dizzy, to touch and to peel.}

Complete the last lesson of chapter 279: I can say to recover, to swell, to feel dizzy, to touch and to peel.
\end{chapteropening}

\section[\texorpdfstring{\te{కోలుకో} (kōlukō)}{కోలుకో (kōlukō)}]{\te{కోలుకో} (kōlukō) — to recover (from illness)}

\label{lesson:TE-C279-koluko}



\begin{quote}
Before the new one: say the Telugu for to cross, then the Telugu for to drive.
\end{quote}

\subsection*{You'll want to know: \te{కోలుకో}}
\textbf{\te{కోలుకో}} — \emph{kōlukō} — \textquotedblleft{}to recover (from illness)\textquotedblright{}.

\begin{grammarlens}[title={What you've built}]
One more word for this chapter.
\end{grammarlens}

\subsection*{Your turn}
\begin{itemize}
\raggedright
  \item \emph{Say it:} \emph{kōlukō}
  \item \emph{Say it:} \emph{kōlukō}, once more
  \item \emph{Say it:} say \emph{kōlukō}
  \item \emph{Recall:} say the Telugu for to cross, then the Telugu for to drive, then say \emph{kōlukō} again
\end{itemize}

\begin{tcolorbox}[breakable,colback=teal!4,colframe=teal!35!black,title={Before you move on}]
What does \te{కోలుకో} mean? (\textquotedblleft{}To recover (from illness)\textquotedblright{}.) Say it once more.
\end{tcolorbox}
\section[\texorpdfstring{\te{ఉబ్బు} (ubbu)}{ఉబ్బు (ubbu)}]{\te{ఉబ్బు} (ubbu) — to swell}

\label{lesson:TE-C279-ubbu}



\begin{quote}
Before the new one: say the Telugu for to drive, then the Telugu for to recover.
\end{quote}

\subsection*{You'll want to know: \te{ఉబ్బు}}
\textbf{\te{ఉబ్బు}} — \emph{ubbu} — \textquotedblleft{}to swell\textquotedblright{}.

\begin{grammarlens}[title={What you've built}]
One more word for this chapter.
\end{grammarlens}

\subsection*{Your turn}
\begin{itemize}
\raggedright
  \item \emph{Say it:} \emph{ubbu}
  \item \emph{Say it:} \emph{ubbu}, once more
  \item \emph{Say it:} say \emph{ubbu}
  \item \emph{Recall:} say the Telugu for to drive, then the Telugu for to recover, then say \emph{ubbu} again
\end{itemize}

\begin{tcolorbox}[breakable,colback=teal!4,colframe=teal!35!black,title={Before you move on}]
What does \te{ఉబ్బు} mean? (\textquotedblleft{}To swell\textquotedblright{}.) Say it once more.
\end{tcolorbox}
\section[\texorpdfstring{\te{కళ్ళు} \te{తిరుగు} (kaḷḷu tirugu)}{కళ్ళు తిరుగు (kaḷḷu tirugu)}]{\te{కళ్ళు} \te{తిరుగు} (kaḷḷu tirugu) — to feel dizzy}

\label{lesson:TE-C279-kallu-tirugu}



\begin{quote}
Before the new one: say the Telugu for to recover, then the Telugu for to swell.
\end{quote}

\subsection*{You'll want to know: \te{కళ్ళు} \te{తిరుగు}}
\textbf{\te{కళ్ళు} \te{తిరుగు}} — \emph{kaḷḷu tirugu} — \textquotedblleft{}to feel dizzy\textquotedblright{}.

\begin{grammarlens}[title={What you've built}]
One more word for this chapter.
\end{grammarlens}

\subsection*{Your turn}
\begin{itemize}
\raggedright
  \item \emph{Say it:} \emph{kaḷḷu tirugu}
  \item \emph{Say it:} \emph{kaḷḷu tirugu}, once more
  \item \emph{Say it:} say \emph{kaḷḷu tirugu}
  \item \emph{Recall:} say the Telugu for to recover, then the Telugu for to swell, then say \emph{kaḷḷu tirugu} again
\end{itemize}

\begin{tcolorbox}[breakable,colback=teal!4,colframe=teal!35!black,title={Before you move on}]
What does \te{కళ్ళు} \te{తిరుగు} mean? (\textquotedblleft{}To feel dizzy\textquotedblright{}.) Say it once more.
\end{tcolorbox}
\section[\texorpdfstring{\te{తాకు} (tāku)}{తాకు (tāku)}]{\te{తాకు} (tāku) — to touch}

\label{lesson:TE-C279-taku}



\begin{quote}
Before the new one: say the Telugu for to swell, then the Telugu for to feel dizzy.
\end{quote}

\subsection*{You'll want to know: \te{తాకు}}
\textbf{\te{తాకు}} — \emph{tāku} — \textquotedblleft{}to touch\textquotedblright{}.

\begin{grammarlens}[title={What you've built}]
One more word for this chapter.
\end{grammarlens}

\subsection*{Your turn}
\begin{itemize}
\raggedright
  \item \emph{Say it:} \emph{tāku}
  \item \emph{Say it:} \emph{tāku}, once more
  \item \emph{Say it:} say \emph{tāku}
  \item \emph{Recall:} say the Telugu for to swell, then the Telugu for to feel dizzy, then say \emph{tāku} again
\end{itemize}

\begin{tcolorbox}[breakable,colback=teal!4,colframe=teal!35!black,title={Before you move on}]
What does \te{తాకు} mean? (\textquotedblleft{}To touch\textquotedblright{}.) Say it once more.
\end{tcolorbox}
\section[\texorpdfstring{\te{ఒలుచు} (olucu)}{ఒలుచు (olucu)}]{\te{ఒలుచు} (olucu) — to peel}

\label{lesson:TE-C279-olucu}



\begin{quote}
Before the new one: say the Telugu for to feel dizzy, then the Telugu for to touch.
\end{quote}

\subsection*{You'll want to know: \te{ఒలుచు}}
\textbf{\te{ఒలుచు}} — \emph{olucu} — \textquotedblleft{}to peel\textquotedblright{}.

\begin{grammarlens}[title={What you've built}]
One more word for this chapter.
\end{grammarlens}

\subsection*{Your turn}
\begin{itemize}
\raggedright
  \item \emph{Say it:} \emph{olucu}
  \item \emph{Say it:} \emph{olucu}, once more
  \item \emph{Say it:} say \emph{olucu}
  \item \emph{Recall:} say the Telugu for to feel dizzy, then the Telugu for to touch, then say \emph{olucu} again
\end{itemize}

\begin{tcolorbox}[breakable,colback=teal!4,colframe=teal!35!black,title={Before you move on}]
What does \te{ఒలుచు} mean? (\textquotedblleft{}To peel\textquotedblright{}.) Say it once more.
\end{tcolorbox}
